\chapter{Introducción}\label{cap:introduccion}

\section{Contexto y motivación}
La movilidad urbana genera volúmenes de datos que exceden las capacidades de
procesamiento tradicional. En este escenario, Apache Spark se ha consolidado como una
tecnología de referencia para analítica distribuida por su modelo de ejecución y su
optimización de consultas \cite{zaharia2012resilient,armbrust2015sparksql,sparkdocs}.
Para la exploración interactiva y la comunicación de resultados, Apache Zeppelin aporta
un formato notebook que integra código, consultas y visualizaciones en un único flujo
de trabajo \cite{zeppelindocs}.

El presente trabajo aplica este stack al conjunto de datos público de la
\textit{New York City Taxi and Limousine Commission} (NYC TLC), centrado en registros
Yellow Taxi en formato Parquet \cite{nyctlc,nyctlcparquet,parquetdocs}. El interés no
es únicamente técnico: se busca demostrar cómo convertir datos masivos de viajes en
evidencia analítica clara para apoyar decisiones sobre patrones de demanda, coste y
comportamiento temporal de la movilidad.

Además, se adopta una arquitectura reproducible en entorno local con Docker, de forma
que cualquier miembro del equipo pueda replicar la ejecución y validar resultados sin
dependencias manuales difíciles de mantener \cite{dockerdocs}.

\section{Objetivos}
\subsection{Objetivo general}
Diseñar e implementar un sistema analítico Big Data reproducible y defendible,
capaz de procesar, limpiar, enriquecer, consultar y visualizar datos masivos de
viajes de taxi en Nueva York, culminando en notebooks de Zeppelin que
concentran todo el análisis técnico y las visualizaciones geográficas.

\subsection{Objetivos específicos}
\begin{itemize}
	\item Cargar y validar ficheros Yellow Taxi en formato Parquet dentro del
	entorno Spark.
	\item Definir y aplicar reglas trazables de calidad del dato directamente
	en los notebooks de Zeppelin.
	\item Generar variables derivadas para análisis temporal, geográfico y económico.
	\item Integrar Spark SQL y DataFrames para consultas analíticas reproducibles.
	\item Enriquecer geográficamente los datos mediante GeoJSON de zonas TLC.
	\item Exportar indicadores por zona a CSV para alimentar mapas coropléticos.
	\item Construir visualizaciones geográficas interactivas con Folium.
\end{itemize}

\section{Alcance}
El alcance funcional del proyecto incluye:
\begin{itemize}
	\item Ejecución local con Docker Compose y servicios Spark + Zeppelin.
	\item Carga, limpieza y enriquecimiento de datos Yellow Taxi (enero 2025).
	\item Análisis temporal y geográfico-económico en notebooks interactivos.
	\item Exportación de indicadores por zona para mapas coropléticos con Folium.
	\item Generación de mapas interactivos que permitan inspección por zona.
\end{itemize}

Quedan fuera del alcance:
\begin{itemize}
	\item Despliegue productivo en clúster empresarial multiusuario.
	\item Orquestación avanzada (por ejemplo, Airflow) y monitorización completa.
	\item Validación con datos de múltiples años de forma exhaustiva en infraestructura
	de alta disponibilidad.
\end{itemize}

Este recorte de alcance es intencional y coherente con el contexto docente: prioriza
calidad técnica, trazabilidad y capacidad de defensa frente a complejidad operativa.

\section{Contribuciones del trabajo}
Las contribuciones principales son las siguientes:
\begin{itemize}
	\item \textbf{Flujo analítico integrado en Zeppelin}: dos notebooks que cubren
	desde la carga hasta la visualización geográfica, sin scripts externos.
	\item \textbf{Criterios de calidad explícitos}: reglas concretas para filtrar registros
	inconsistentes y justificar su impacto.
	\item \textbf{Enriquecimiento geográfico completo}: uso de GeoJSON y Folium para
	representar indicadores por zona TLC con mapas coropléticos interactivos
	\cite{foliumdocs,geojsonspec}.
	\item \textbf{Exportación de indicadores por zona}: generación de CSVs de métricas
	agregadas por origen y destino para alimentar visualizaciones externas.
	\item \textbf{Narrativa defendible}: documentación técnica que conecta decisiones,
	implementación y resultados sin introducir métricas inventadas.
\end{itemize}

\section{Estructura del documento}
El documento se organiza en siete capítulos principales. Tras esta introducción,
el Capítulo 2 describe el estudio previo y la gestión del trabajo. El Capítulo 3
formaliza los requisitos y el análisis del problema. El Capítulo 4 presenta el
diseño de la arquitectura y del flujo analítico. El Capítulo 5 detalla la implementación
real del repositorio y los notebooks de Zeppelin. El Capítulo 6 recoge
las pruebas de validación funcional y técnica. Por último, el capítulo de
conclusiones resume logros, limitaciones y líneas de trabajo futuro.

